\chapter{Conclusions}\label{ch:conclusions}
\thispagestyle{pagestyle}

\section{Summary of contributions}
This thesis studied a NISQ-oriented structural optimization of FRQI state preparation for a fixed QHED edge-detection stage, in a complete end-to-end pipeline.

The following claims can be made based on \Cref{ch:results}:

\begin{itemize}
  \item \textbf{Ideal FRQI encoding:} SSIM metric reports parity between pre- and post- FRQI encoding for all image sizes.
  \item \textbf{QHED vs. Sobel:} QHED diverges from classical Sobel as image size grows, consistent with prior literature, not because one is better than the other, but rather due to their different approach on reporting.
  \item \textbf{Structural resources:} Replacing the naive \texttt{noancilla} MCX decomposition with a v-chain ancilla pattern cuts transpiled CX significantly at larger images, at the cost of a modest number of qubits.
  \item \textbf{Ideal and noisy image quality:} Under depolarizing noise, QHED is not having improvements when comparing FRQI state preparation naive versus v-chain. This means that lower CX depth does not give better edge detection quality, which remains true for all noise scales.
  \item \textbf{Readout mitigation:} An improvement can be observed when applying an inversion of the calibration matrix to correct the readout bias.
\end{itemize}

\section{Future work}

Promising extensions to this thesis include:

\begin{itemize}
    \item running the experiments on real IBM hardware, to understand the differences between simulated and real environments,
    \item completing the $16\times 16$ v-chain density-matrix experiments, which were not possible due to memory constraints,
    \item zero-noise extrapolation and richer mitigation stacks beyond the size of $2\times 2$ readout inversion,
    \item comparing additional \texttt{MCXGate} modes from recent papers,
    \item hybrid edge schemes that keep shallow quantum layers while adopting cheaper FRQI preparation.
\end{itemize}